一辆行驶的汽车会不断消耗汽油，飞行中的飞机或火箭也有同样的特点：运动过程中，
研究对象的质量不断变化。处理这类问题时，必须同时计入速度变化和质量变化带来的
动量变化。变质量并不意味着动量定理失效；真正需要重新说明的是研究系统的边界——哪一
部分物质在当前时刻属于研究对象，哪些物质正从边界流入或流出，以及它们穿过边界前后的
速度分别是多少。

先只对主体动量作代数展开，不把它等同于外力冲量：
\begin{equation}
 \Delta\vec p=\Delta(m\vec v)
 =\Delta m\,\vec v+m\,\Delta\vec v+\Delta m\,\Delta\vec v.
\end{equation}
略去二阶小量并取连续极限，得到严格的微分关系
\[
 \dd\vec p=\vec v\,\dd m+m\,\dd\vec v.
\]
应用这一关系时还要注意：新进入系统或刚离开系统的物质，未必具有研究对象本身的速度。
这些物质所带入或带走的动量，正是变质量问题中容易遗漏的一项。

\subsubsection*{变质量系统的边界与动量通量}

火箭、落链和接雨小车的研究对象会随时间增减质量，不能直接把 $\vec F=m\vec a$ 中的
$m$ 换成 $m(t)$。动量定理适用于同一批物质；若研究对象在两个时刻包含的物质不同，
就要把跨越边界所带的动量单独记入。

设主体质量为 $m$、速度为 $\vec v$，时间 $\dd t$ 内吸收质量 $\dd m_{\rm in}>0$，
排出质量 $\dd m_{\rm out}>0$。流入、流出物质在同一个惯性系中的速度分别为
$\vec u_{\rm in},\vec u_{\rm out}$，主体质量改变
$\dd m=\dd m_{\rm in}-\dd m_{\rm out}$。把主体及这次交换的物质选为同一批物质：
\[
 m\vec v+\vec u_{\rm in}\dd m_{\rm in}+\vec F_{\rm ext}\dd t
 =(m+\dd m)(\vec v+\dd\vec v)+\vec u_{\rm out}\dd m_{\rm out}.
\]
这里 $\vec F_{\rm ext}$ 是连续极限下作用于主体的外力；微小交换质量所受有限外力的冲量
为二阶小量。展开、略去二阶项并除以 $\dd t$，得到
\begin{equation}\label{eq:mass-flux-balance}
 \frac{\dd(m\vec v)}{\dd t}
 =\vec F_{\rm ext}+\dot m_{\rm in}\vec u_{\rm in}
 -\dot m_{\rm out}\vec u_{\rm out}.
\end{equation}
左侧是主体动量增加率，右侧依次为外力、流入动量率、流出动量率。展开左侧乘积，得
\begin{equation}\label{eq:mass-acceleration}
 m\frac{\dd\vec v}{\dd t}
 =\vec F_{\rm ext}+\dot m_{\rm in}(\vec u_{\rm in}-\vec v)
 -\dot m_{\rm out}(\vec u_{\rm out}-\vec v).
\end{equation}
只有一股交换流时，也可用有符号的 $\dot m$（流入为正、流出为负）和相对主体速度
$\vec u_{\rm rel}=\vec u-\vec v$ 写成 $m\dot{\vec v}=\vec F_{\rm ext}+\vec u_{\rm rel}\dot m$。
不能把这个相对速度与前一式的惯性系绝对速度混用。动量通量的量纲为
$\mathrm{kg/s}\times\mathrm{m/s}=\mathrm N$，与外力一致。

理想火箭在无外力的一维运动中，以相对火箭大小为 $u$ 的速度向后喷气，满足
\begin{equation}
 m\,\dd v=-u\,\dd m,
 \qquad
 v-v_0=u\ln\frac{m_0}{m}.
\end{equation}
速度增量取决于质量比而不是燃料质量的简单差值。若有重力或阻力，还要把它们在燃烧时间内的
冲量计入；若喷气速度相对地面给出，必须先换成相对火箭速度。
火箭自身不是固定物质系统，即使无外力，其自身动量也不必守恒；守恒的是火箭加燃气的总动量。

链条落到秤盘、雨水落入车厢等问题还包含“捕获并加速新质量”的作用。流入物质在进入系统前
已有速度，主体必须提供冲量把它变为共同速度，于是支持力或阻力中出现 $\dot m\Delta v$ 项。
柔链问题还应区分移动段、静止段和弯折区，弯折区虽很短，却承担改变动量方向的有限合力。
检验答案时可看两个极限：质量流率趋于零时应回到定质量问题；相对流速趋于零时，动量通量项
也应消失。最后一种情形只说明无需额外改变主体速度，主体总动量仍可随它包含的质量改变。


\begin{origexamplebox}{例 3.24}
在雨中用恒力 $F$ 推一个质量为 $M$、开口截面积为 $S$ 的容器。若一个静止在地面的
同样容器中，单位时间内水位上升 $\Delta h$，不考虑雨水对容器前方的阻力，水的密度
为 $\rho$，求在 $3$ 个单位时间后容器的速度。
容器初始静止；一个单位时间记为时长 $\tau$，所求时刻就是 $t=3\tau$。
\end{origexamplebox}

\begin{xesolution}
设水平方向为推力方向。由时长 $\tau$ 内落入的水体积 $S\Delta h$，求得质量流率
\[
 q=\frac{\rho S\Delta h}{\tau}.
\]
雨滴落入前没有水平速度，所以它们带入系统的水平动量为零。经过时间 $t$，容器及其中
雨水的总质量为 $M+qt$；恒力的冲量为 $Ft$。由水平动量定理，
\[
 (M+qt)v=Ft,
\]
故
\[
 v(t)=\frac{Ft}{M+(\rho S\Delta h/\tau)t}.
\]
图~\ref{ch07:fig:soldiagThreeTwentyFour}中，雨水只增加质量，不带入水平动量；外力使系统水平动量按 $Ft$ 增长。
\soldiagThreeTwentyFour
代入 $t=3\tau$，得到
\[
 \boxed{v=\frac{3F\tau}{M+3\rho S\Delta h}}.
\]
若初速度为 $v_i$，应保留初始动量，把分子改为 $Mv_i+Ft$。单位检验：$q$ 是质量/时间，
$Ft$ 是冲量，故结果为速度。不能直接写 $F=(M+qt)a$，因为还需要将新进入的雨水加速；
展开正确的动量式可得 $(M+qt)\dot v=F-qv$。当 $q\to0$ 时回到 $v=Ft/M$。
\end{xesolution}

\begin{origexamplebox}{例 3.25}
一根质量为 $m$、长度为 $l$ 的柔软绳子放在水平桌面上。用手捏住绳的一端，以恒定速度
$v$ 竖直向上提绳。当这一端离桌面的高度为 $x<l$ 时，求手的拉力大小。
\end{origexamplebox}

\begin{xesolution}
绳的线密度为
\[
 \lambda=\frac{m}{l}.
\]
高度为 $x$ 时，竖直运动绳段的质量是 $\lambda x$，其速度保持为 $v$。在时间
$\dd t$ 内，又有长度 $v\dd t$、质量 $\lambda v\dd t$ 的绳元由静止转为以速度 $v$
向上运动，因此单位时间增加的向上动量为 $\lambda v^2$。

对竖直运动部分列动量方程：
\[
 F-\lambda xg=\lambda v^2.
\]
图~\ref{ch07:fig:soldiagThreeTwentyFive}中，提力既承担运动绳段的重量，也为新进入运动段的绳元提供动量。
\soldiagThreeTwentyFive
所以
\[
 \boxed{F=\lambda xg+\lambda v^2
 =\frac{m}{l}\left(xg+v^2\right)}.
\]
第一项承担已提起绳段的重量，第二项用来连续加速刚离开桌面的绳元。该结论采用理想柔绳、
转折区域很小且绳元离桌后立即获得速度 $v$ 的模型。
\end{xesolution}

\Needspace{0.38\textheight}
\begin{origexamplebox}{例 3.26}
如图，一根质量为 $m$、长度为 $l$ 的均匀柔绳竖直悬挂，下端刚好与台秤秤盘接触。
释放绳的上端，求从释放开始台秤示数随时间的变化，并求台秤的最大示数。
图~\ref{ch07:fig:probdiagThreeTwentySix}中，柔绳上端释放后下落，绳元依次落到秤盘上并停止。
\probdiagThreeTwentySix
\end{origexamplebox}

\begin{xesolution}
设绳的线密度为 $\lambda=m/l$，释放后经过时间 $t$，绳的上端下落距离为 $x$。理想模型
中尚未落到秤盘上的绳段作自由落体，因而
\[
 x=\frac12gt^2,\qquad v=gt,\qquad v^2=2gx.
\]
秤盘的受力来源有两部分：
\[
 \underbrace{\lambda xg}_{\text{已停在盘上的绳重}}
 \quad\text{和}\quad
 \underbrace{\lambda v^2}_{\text{使下落绳元停止的动量流}}.
\]
图~\ref{ch07:fig:soldiagThreeTwentySix}中，秤盘既承受已落绳段的重量，也连续吸收下落绳元的动量。
\soldiagThreeTwentySix
所以在整根绳尚未全部落下时，
\[
 N=\lambda xg+\lambda v^2
 =3\lambda gx
 =\frac32\lambda g^2t^2.
\]
上端落到秤盘所需时间为
\[
 t_0=\sqrt{\frac{2l}{g}}.
\]
因此理想化的示数为
\[
 \boxed{
 N(t)=
 \begin{cases}
 \dfrac{3m g^2}{2l}\,t^2,
   &0\le t<\sqrt{\dfrac{2l}{g}},\\[6pt]
 mg,&t>\sqrt{\dfrac{2l}{g}}.
 \end{cases}}
\]
在全部落下前的一瞬间，
\[
 \boxed{N_{\max}=3mg}.
\]
随后动量流消失，稳定示数为 $mg$。理想模型在 $t=t_0$ 处给出突变；真实柔绳因弯曲、
碰撞持续时间和秤的响应时间有限，峰值会被平滑。
\end{xesolution}

\begin{origexamplebox}{例 3.27}
一个截面积为 $S$ 的物体在密度均匀、密度为 $\rho$ 的介质中竖直下落。物体会把前方
遇到的介质全部吸收，忽略其他阻力。写出物体速度 $v$ 和质量 $m$ 随时间变化的方程，
并定性说明其运动规律。
\end{origexamplebox}

\begin{xesolution}
取竖直向下为正。时间 $\dd t$ 内物体扫过体积 $Sv\dd t$，吸收的静止介质质量为
\[
 \dd m=\rho Sv\,\dd t,
 \qquad
 \boxed{\dot m=\rho Sv}.
\]
被吸收介质在地面系中初始动量为零，因此“物体加上刚吸收的介质”的动量变化只由重力
产生：
\[
 \frac{\dd(mv)}{\dd t}=mg.
\]
展开后得到
\[
 \boxed{m\dot v=mg-\rho Sv^2}.
\]
\Needspace{10\baselineskip}
图~\ref{ch07:fig:soldiagThreeTwentySeven}中，静止介质被吸收后要加速到 $v$，故出现随 $v^2$ 增长的动量流项。
\soldiagThreeTwentySeven
右端第二项不是通常意义上的阻力，而是把静止介质加速到速度 $v$ 所需要的动量流。若
初始时物体由静止释放，质量为 $m_0$，令 $k=\rho S$，消去时间：
\[
 \frac{\dd(v^2)}{\dd m}+\frac{2}{m}v^2=\frac{2g}{k}.
\]
用积分因子 $m^2$ 将左侧合为乘积的导数：
\[
 \frac{\dd}{\dd m}(m^2v^2)=\frac{2g}{k}m^2,
 \qquad m^2v^2=\frac{2g}{3k}(m^3-m_0^3).
\]
除以 $m^2$，得
\[
 v^2=\frac{2g}{3k}
 \left(m-\frac{m_0^3}{m^2}\right).
\]
因此 $m$、$v$ 均随时间增大，但加速度
\[
 \dot v=g-\frac{k}{m}v^2=\frac g3+\frac{2g}{3}\left(\frac{m_0}{m}\right)^3
\]
逐渐减小，并趋近于 $g/3$；速度并不趋于常量。
若初速度不为零，质量式与速度式仍成立，但单调性须重新检查；初始条件
$\rho S v_i^2>m_0g$ 会使物体先减速，不能把从静止出发的结论直接推广。
\end{xesolution}

\Needspace{0.30\textheight}
\begin{origexamplebox}{例 3.28}
一个球形小水滴在均匀的雾中下落。设雾的质量密度均匀，水滴会吸收所遇到的全部小水滴，
忽略空气阻力，重力加速度 $g$ 为常量。写出大水滴的质量、速度和半径随时间变化的方程，
并定性说明其运动规律。
\end{origexamplebox}

\begin{xesolution}
设雾的质量密度为 $\rho_f$，水的密度为 $\rho_w$，大水滴半径为 $r$。则
\[
 m=\frac43\pi\rho_w r^3,
 \qquad
 \dot m=\rho_f\pi r^2v.
\]
由第一式求导并与第二式比较：
\[
 4\pi\rho_w r^2\dot r=\rho_f\pi r^2v,
 \qquad
 \boxed{\dot r=\frac{\rho_f}{4\rho_w}v}.
\]
雾滴在被吸收前近似静止，故动量方程仍为
\[
 \frac{\dd(mv)}{\dd t}=mg,
\]
即
\[
 \boxed{\dot v
 =g-\frac{\dot m}{m}v
 =g-\frac{3\rho_f}{4\rho_w r}v^2}.
\]
\Needspace{10\baselineskip}
图~\ref{ch07:fig:soldiagThreeTwentyEight}中，半径、迎风面积和质量共同增长，速度必须与它们联立求解。
\soldiagThreeTwentyEight
若水滴由静止开始、初始半径为 $r_0$，令
$\alpha=\rho_f/(4\rho_w)$，令 $y=v^2$，由 $\dot r=\alpha v$ 消去时间，得
\[
 \frac{\dd y}{\dd r}+\frac{6}{r}y=\frac{2g}{\alpha}.
\]
乘以积分因子 $r^6$ 并积分：
\[
 \frac{\dd}{\dd r}(r^6y)=\frac{2g}{\alpha}r^6,
 \qquad r^6y=\frac{2g}{7\alpha}(r^7-r_0^7).
\]
因此
\[
 v^2=\frac{2g}{7\alpha}
 \left(r-\frac{r_0^7}{r^6}\right).
\]
代回原方程可得 $\dot v=g/7+6g(r_0/r)^7/7>0$。所以对这一静止初值，
$r$、$m$ 和 $v$ 都不断增大；随着吸收过程进行，加速度由初始的 $g$ 逐渐减小，
并趋近于 $g/7$。半径增大使迎风面积增大，吸收增强；质量增长又改变速度，三者必须联立
考虑。
若初始速度为 $v_i$，当 $3\alpha v_i^2/r_0>g$ 时会先减速。这里用 $\alpha$ 表示密度比，
不与加速度符号 $a$ 混用；$gr/\alpha$ 的量纲为速度平方。
\end{xesolution}

火箭依靠从尾部喷出燃料燃烧产生的高温高压气体而获得巨大推力。实际升空过程十分复杂，
为简化分析，下面只研究火箭沿直线飞行的理想模型。忽略重力和其他外力，设火箭喷出的
燃气相对火箭向后速度大小为 $u$。在极短时间 $\Delta t$ 内，喷气前火箭的质量和速度为 $M,v$；
喷气后为 $M+\Delta M,v+\Delta v$，其中 $\Delta M<0$。喷出燃气的质量为
$-\Delta M$，在地面参考系中的速度为 $v-u$。
图~\ref{ch07:fig:rocketMicroDiagram}中，在地面系中，喷出燃气的速度是 $v-u$；$\Delta M<0$。
\rocketMicroDiagram
由喷气前后的总动量守恒，
\[
 Mv=(M+\Delta M)(v+\Delta v)-\Delta M(v-u).
\]
展开并消去 $Mv$：
\[
 M\Delta v+u\Delta M+\Delta M\Delta v=0.
\]
略去二阶小量 $\Delta M\Delta v$，得到
\begin{equation}
 \label{eq:rocket-micro}
 \Delta v=-u\frac{\Delta M}{M}.
\end{equation}
因为 $\Delta M<0$，所以 $\Delta v>0$：火箭向后喷出质量，自身向前加速。在连续极限下，
式 \eqref{eq:rocket-micro} 可写成
\[
 \dd v=-u\frac{\dd M}{M}.
\]

\Needspace{0.34\textheight}
\begin{origexamplebox}{例 3.29}
一枚火箭在某时刻的质量为 $M$、速度为 $v$，燃气相对火箭的喷出速度大小为 $u$。
忽略重力，请根据式 \eqref{eq:rocket-micro}，在下面坐标系中定性画出此后火箭速度与质量的关系。
图~\ref{ch07:fig:probdiagThreeTwentyNine}中，曲线应从当前状态 $(M,v)$ 出发，并标明燃烧过程的方向。
\probdiagThreeTwentyNine
\end{origexamplebox}

\begin{xesolution}
设此后某时刻火箭的剩余质量为 $\mathcal M$、速度为 $V$。由
\[
 \dd V=-u\frac{\dd\mathcal M}{\mathcal M}
\]
从当前状态 $(M,v)$ 积分到 $(\mathcal M,V)$：
\[
 V-v=-u\int_M^{\mathcal M}\frac{\dd m}{m}
 =u\ln\frac{M}{\mathcal M}.
\]
因此
\[
 \boxed{V(\mathcal M)=v+u\ln\frac{M}{\mathcal M}},
 \qquad 0<\mathcal M\le M.
\]
图~\ref{ch07:fig:soldiagThreeTwentyNine}中，横轴质量减小时，火箭沿对数曲线向左上方运动。
\soldiagThreeTwentyNine
曲线经过 $(M,v)$，并满足
\[
 \frac{\dd V}{\dd\mathcal M}=-\frac{u}{\mathcal M}<0,
 \qquad
 \frac{\dd^2V}{\dd\mathcal M^2}
 =\frac{u}{\mathcal M^2}>0.
\]
所以按“横轴质量、纵轴速度”作图时，曲线单调下降，且斜率随质量增大逐渐变平（曲线
开口向上）；实际燃烧过程的方向是沿曲线向左上方运动，即质量减小、速度增大。
$\mathcal M\to0$ 时速度发散只是理想模型的数学
外推，真实火箭不可能把结构质量全部喷掉。
\end{xesolution}

\Needspace{0.30\textheight}
\begin{origexamplebox}{例 3.30}
一枚单级火箭的初始质量为 $M_0$，燃料燃烧完毕后的质量为 $M$。燃气相对火箭的喷出速度
大小恒为 $u$，火箭由静止开始运动，忽略重力和其他阻力。证明火箭的末速度为
\[
 v=u\ln\frac{M_0}{M}.
\]
\end{origexamplebox}

\begin{xesolution}
由式 \eqref{eq:rocket-micro} 的连续形式，
\[
 \dd v=-u\frac{\dd M}{M}.
\]
火箭由静止开始，质量由 $M_0$ 减少到 $M$。对整个燃烧过程积分：
\[
 \int_0^v\dd v
 =-u\int_{M_0}^{M}\frac{\dd M}{M}.
\]
图~\ref{ch07:fig:soldiagThreeThirty}中，对 $\dd v=-u\,\dd M/M$ 积分，速度增量等于 $u/M$ 曲线下对应的面积。
\soldiagThreeThirty
于是
\[
 v=-u\left(\ln M-\ln M_0\right),
\]
即
\[
 \boxed{v=u\ln\frac{M_0}{M}}.
\]
这就是理想火箭方程。若定义质量比 $R=M_0/M$，则速度增量为
$\Delta v=u\ln R$：提高喷气速度 $u$ 能线性提高速度增量，而增大质量比的收益只有
对数增长。证明成立的条件是喷气相对速度恒定，并忽略重力、空气阻力等外力。
\end{xesolution}
